\section*{Blow-up classification (author source pp.38--39)}
\begin{originalresult}{Definition 6.4}
We say that a gadget $g:X\times Y\to\bits$ is a blow-up of a gadget $h:Z\times W\to\bits$ if there are decompositions $X=\bigsqcup_{z\in Z}X_z$ and $Y=\bigsqcup_{w\in W}Y_w$, with $X_z,Y_w\ne\varnothing$, such that all $(x_z,y_w)\in X_z\times Y_w$ satisfy $g(x_z,y_w)=h(z,w)$.
\end{originalresult}
For example, $g$ is a blow-up of XOR if (up to rearrangement) it has a communication matrix of the following form
\[\begin{bmatrix}1&\cdots&1&0&\cdots&0\\\vdots&\ddots&\vdots&\vdots&\ddots&\vdots\\1&\cdots&1&0&\cdots&0\\0&\cdots&0&1&\cdots&1\\\vdots&\ddots&\vdots&\vdots&\ddots&\vdots\\0&\cdots&0&1&\cdots&1\end{bmatrix}.\]
If $g$ is a blow-up of $h$, then $f\comp g$ and $f\comp h$ have the same communication complexity in all models. Indeed, on the one hand, $h$ is a restriction of $g$, and so a protocol for $f\comp g$ can be used to solve $f\comp h$; and on the other hand, by replacing $x\in X_z$ by $z$ and $y\in Y_w$ by $w$, we can use a protocol for $f\comp h$ to solve $f\comp g$.
\begin{originalresult}{Corollary 6.5}
If $g:X\times Y\to\bits$ satisfies both $\OR\nleq g$ and $\AND\nleq g$ then either $g$ is constant, or it is a blow-up of XOR.
\end{originalresult}
\begin{proof}
Let $X=X_0\sqcup X_1\sqcup\cdots\sqcup X_k$ and $Y=Y_0\sqcup Y_1\sqcup\cdots\sqcup Y_k$ be the decomposition provided by Lemma 6.3. We can assume that the $X_i,Y_i\ne\varnothing$, except that possibly $X_0=Y_0=\varnothing$.
If $X_0\ne\varnothing$ and $k>0$ then $\AND\le g$. Indeed, if we restrict the matrix of $g$ so that it contains one row each from $X_0,X_1$ and one column each from $Y_0,Y_1$ then we get
\[\begin{bmatrix}0&0\\0&1\end{bmatrix}\]
which is AND. Consequently, either $k=0$ or $X_0=Y_0=\varnothing$.
If $k=0$ then $g\equiv0$. Now suppose that $X_0=Y_0=\varnothing$. If $k\ge3$ and we restrict the communication matrix of $g$ so that it contains one row from each $X_1,X_2,X_3$ and one column from each $Y_1,Y_2,Y_3$ then we get
\[\begin{bmatrix}\mathbf1&0&\mathbf0\\\mathbf0&1&\mathbf0\\0&0&1\end{bmatrix}.\]
Again, we get a restriction to AND (the entries in bold text correspond to the AND matrix). Hence, $k\le2$. If $k=1$ then $g\equiv1$, and otherwise $g$ is a blow-up of XOR.
\end{proof}
\begin{originalresult}{Corollary 6.6}
If $g:X\times Y\to\bits$ satisfies both $\OR\nleq g$ and $\XOR\nleq g$ then either $g$ is constant, or it is a blow-up of AND. Similarly, if $\AND\nleq g$ and $\XOR\nleq g$ then either $g$ is constant, or it is a blow-up of OR.
\end{originalresult}
\begin{proof}
We can assume that the $X_i,Y_i\ne\varnothing$, except that possibly $X_0=Y_0=\varnothing$. We start by proving the first claim, and so we assume that $\OR,\XOR\nleq g$.
If $k\ge2$ then $\XOR\le g$. Indeed, if we restrict the matrix of $g$ so that it contains one row from each $X_1,X_2$ and one column from each $Y_1,Y_2$ then we obtain
\[\begin{bmatrix}1&0\\0&1\end{bmatrix}\]
which is XOR. If $k=0$ then $g\equiv0$. If $k=1$ and $X_0=Y_0=\varnothing$ then $g\equiv1$. Finally, if $k=1$ and $X_0,Y_0\ne\varnothing$ then $g$ is a blow-up of AND.
Suppose that $\AND,\XOR\nleq g$. Then $\OR,\XOR\nleq\bar g$, where $\bar g$ is the negation of $g$. Hence, either $\bar g$ is constant or it is a blow-up of AND. Thus, either $g$ is constant or it is a blow-up of OR.
\end{proof}
